\chapter{约束极值}
\label{chap:constrained-extrema}

求函数在曲线、曲面或若干方程限定的集合上的极值时，允许的微小位移不再是
整个空间中的任意方向，而只能沿约束集的切方向。Lagrange 乘子条件是这个
事实的线性代数表述：目标函数的微分在所有可行切向上为零，因而它属于约束
微分所张成的空间。

乘子方程只是必要条件。要保证它确实覆盖所有局部极值，约束的梯度必须满足
正则性；要判定候选点，还需把二阶型限制到切空间。对不等式约束，只有在
极值点取等号的约束参与一阶条件，并出现非负乘子和互补松弛关系。

\section{等式约束的切空间}

设 $U\subset\R^n$ 开，$h=(h_1,\ldots,h_m):U\to\R^m$ 属于 $C^1$，
并考虑
\[
 M=\{x\in U:h(x)=0\}.
\]
若 $Dh(a)$ 满射，即梯度
$\nabla h_1(a),\ldots,\nabla h_m(a)$ 线性无关，则 $a$ 是约束的正则点。
由正则水平集定理，$M$ 在 $a$ 附近是维数 $n-m$ 的局部图像，并且
\[
 T_aM=\ker Dh(a)
 =\{v:\langle\nabla h_j(a),v\rangle=0, j=1,\ldots,m\}.
\]

下面的线性代数事实把“在切空间上为零”变成乘子表示。

\begin{lemma}[消去核的线性泛函]
\label{lem:functional-factor-through-surjection}
设 $A:\R^n\to\R^m$ 满射，$\ell\in(\R^n)^*$。若
$\ell|_{\ker A}=0$，则存在唯一的 $\lambda\in(\R^m)^*$，使
\[
 \ell=\lambda\circ A.
\]
\end{lemma}

\begin{proof}
对 $y\in\R^m$ 选取任意 $x$ 使 $Ax=y$，定义 $\lambda(y)=\ell(x)$。
若 $Ax=Ax'$，则 $x-x'\in\ker A$，故 $\ell(x)=\ell(x')$；所以定义与
原像选择无关。线性性显然，且 $\ell=\lambda\circ A$。满射性还保证
$\lambda$ 唯一。
\end{proof}

\section{Lagrange 乘子定理}

\begin{theorem}[等式约束的 Lagrange 乘子定理]
\label{thm:lagrange-multipliers-regular-equality}
设 $f,h_1,\ldots,h_m$ 在 $a$ 的邻域属于 $C^1$，$h(a)=0$，且
$Dh(a)$ 满射。若 $a$ 是 $f$ 在 $M=h^{-1}(0)$ 上的局部极值点，则存在
唯一的乘子 $\lambda_1,\ldots,\lambda_m$，使
\begin{equation}
 \nabla f(a)+\sum_{j=1}^m\lambda_j\nabla h_j(a)=0.
 \label{eq:lagrange-stationarity-equality}
\end{equation}
\end{theorem}

\begin{proof}
任取 $v\in T_aM$。正则水平集的局部参数化保证存在 $M$ 中的 $C^1$ 曲线
$\gamma$，满足 $\gamma(0)=a$、$\gamma'(0)=v$。一元函数
$f(\gamma(t))$ 在 $0$ 取得局部极值，所以
$Df(a)v=0$。因此 $Df(a)$ 在 $\ker Dh(a)$ 上为零。

由引理 \ref{lem:functional-factor-through-surjection}，存在唯一
$\mu\in(\R^m)^*$ 使 $Df(a)=\mu\circ Dh(a)$。把
$\mu(y)=-\sum_j\lambda_jy_j$ 写入标准坐标，便得到
\eqref{eq:lagrange-stationarity-equality}。
\end{proof}

将
\[
 L(x,\lambda)=f(x)+\sum_{j=1}^m\lambda_jh_j(x)
\]
称为 Lagrange 函数，则候选点满足
\[
 \nabla_xL(a,\lambda)=0,\qquad h(a)=0.
\]
这些方程需要连同原约束一起求解。解出的点只是候选点；极值类型仍要由
比较、紧致性或二阶条件判定。

\begin{example}[球面上的线性函数]
在 $\norm x=1$ 上求 $f(x)=\langle c,x\rangle$ 的极值。取
$h(x)=\norm x^2-1$，乘子方程为
\[
 c+2\lambda x=0.
\]
若 $c\ne0$，候选点只有 $x=\pm c/\norm c$。Cauchy--Schwarz 不等式
进一步确认它们分别给出全局最大值 $\norm c$ 和最小值 $-\norm c$。
乘子法找到候选点，最后的全局判定来自球面的紧致性和具体不等式。
\end{example}

正则性不能省略。例如约束 $h(x)=x^2=0$ 的可行集只有原点；原点当然是
$f(x)=x$ 在该集合上的极大点和极小点，但 $h'(0)=0$，方程
$f'(0)+\lambda h'(0)=0$ 无解。此处不是极值不存在，而是退化的约束函数
没有正确表示可行集的一阶几何。同一集合若改由 $h(x)=x=0$ 表示，乘子
条件便恢复正常。

\section{二阶条件}

设 $f,h$ 属于 $C^2$，$a$ 是正则约束点，$\lambda$ 满足一阶乘子方程。
沿 $M$ 中曲线 $\gamma$ 二次求导，可得
\[
 (f\circ\gamma)''(0)
 =D^2f(a)[v,v]+Df(a)\gamma''(0),
\]
而由 $h_j(\gamma(t))=0$，
\[
 D^2h_j(a)[v,v]+Dh_j(a)\gamma''(0)=0.
\]
用一阶乘子方程消去 $\gamma''(0)$，得到
\begin{equation}
 (f\circ\gamma)''(0)=D^2_{xx}L(a,\lambda)[v,v]
 \qquad(v=\gamma'(0)\in T_aM).
 \label{eq:constrained-second-variation}
\end{equation}

\begin{theorem}[等式约束的二阶条件]
\label{thm:constrained-second-order-test}
在上述假设下：
\begin{enumerate}[label=\textup{(\roman*)}]
 \item 若 $a$ 是约束局部极小点，则
 $D^2_{xx}L(a,\lambda)[v,v]\ge0$ 对每个 $v\in T_aM$ 成立；约束局部
 极大点满足反向不等式。
 \item 若该二次型在 $T_aM\setminus\{0\}$ 上正定，则 $a$ 是严格约束
 局部极小点；若负定，则为严格约束局部极大点。
\end{enumerate}
\end{theorem}

\begin{proof}
必要性由 \eqref{eq:constrained-second-variation} 和一元二阶必要条件得到。
证明充分性时，取正则水平集的 $C^2$ 局部参数化
$\psi:V\subset\R^{n-m}\to M$，令 $\psi(0)=a$。函数
$F=f\circ\psi$ 在零点的一阶微分为零，其 Hessian 满足
\[
 D^2F(0)[w,w]
 =D^2_{xx}L(a,\lambda)[D\psi(0)w,D\psi(0)w].
\]
$D\psi(0)$ 是从 $\R^{n-m}$ 到 $T_aM$ 的线性同构。切空间上正定因而
推出 $D^2F(0)$ 正定，通常的 Hessian 判别说明 $0$ 是 $F$ 的严格局部
极小点，也就是 $a$ 是约束严格局部极小点。负定情形相同。
\end{proof}

应当限制的是 $D^2_{xx}L$，而不是只限制 $D^2f$。曲面自身的弯曲通过
$\sum_j\lambda_jD^2h_j$ 进入二阶变化。也无须检查整个 $\R^n$ 上的正定性；
法向方向不是一阶可行方向。

\begin{example}
在圆周 $x^2+y^2=1$ 上考察 $f(x,y)=x^2+2y^2$。候选点为
$(\pm1,0)$、$(0,\pm1)$。切空间分别是一维的。把 Lagrange Hessian 限制
到相应切线，可得前两点为严格最小，后两点为严格最大。直接代入圆周参数
也得到同一结论，但切空间方法能直接推广到多重约束。
\end{example}

\section{不等式约束与 KKT 条件}

考虑
\[
 h_j(x)=0\quad(j=1,\ldots,m),
 \qquad
 g_i(x)\le0\quad(i=1,\ldots,r).
\]
在可行点 $a$，满足 $g_i(a)=0$ 的指标组成活跃集 $I(a)$。非活跃约束在
$a$ 附近仍严格小于零，不限制充分小的扰动。

一种常用的约束资格条件是 LICQ：
\[
 \{\nabla h_j(a):1\le j\le m\}
 \cup\{\nabla g_i(a):i\in I(a)\}
\]
线性无关。它既保证等式约束正则，也排除活跃不等式的一阶描述相互退化。

\begin{theorem}[Karush--Kuhn--Tucker 必要条件]
\label{thm:kkt-necessary-licq}
设 $f,h_j,g_i$ 在 $a$ 的邻域属于 $C^1$，$a$ 是上述约束下的局部极小点，
且 LICQ 在 $a$ 成立。则存在 $\lambda_j\in\R$ 和 $\mu_i\ge0$，使
\begin{align}
 &\nabla f(a)+\sum_{j=1}^m\lambda_j\nabla h_j(a)
   +\sum_{i=1}^r\mu_i\nabla g_i(a)=0,
 \label{eq:kkt-stationarity}\\
 &h_j(a)=0,\qquad g_i(a)\le0,
 \label{eq:kkt-primal}\\
 &\mu_i g_i(a)=0\qquad(i=1,\ldots,r).
 \label{eq:kkt-complementarity}
\end{align}
\end{theorem}

\begin{proof}
对非活跃指标先置 $\mu_i=0$。LICQ 允许把
$h_1,\ldots,h_m$ 和活跃的 $g_i$ 补成一个局部坐标系。具体地，选取
若干线性坐标 $z$，使
\[
 \Phi(x)=(h(x),g_{I(a)}(x),z(x))
\]
在 $a$ 的导数可逆；逆函数定理说明 $\Phi$ 是局部坐标变换。把 $a$ 移到
原点，并以 $(u,v,w)=\Phi(x)$ 表示。局部可行集由
\[
 u=0,\qquad v_i\le0
\]
描述，$w$ 可以在两个方向自由变化。

令 $\widetilde f=f\circ\Phi^{-1}$。沿 $w$ 的正负方向都可行，局部极小性
给出 $D_w\widetilde f(0)=0$。沿 $-e_i$ 的 $v$ 方向可行，故
\[
 -\partial_{v_i}\widetilde f(0)\ge0.
\]
取 $\mu_i=-\partial_{v_i}\widetilde f(0)\ge0$，再用自由的 $u$ 分量定义
$\lambda_j=-\partial_{u_j}\widetilde f(0)$。把坐标形式拉回原变量，正好得到
\eqref{eq:kkt-stationarity}。非活跃约束的乘子为零，因此
\eqref{eq:kkt-complementarity} 也成立。
\end{proof}

乘子 $\mu_i$ 的非负性来自单侧可行方向；互补松弛说明只有活跃约束可能对
驻点方程施加法向力。若把约束统一写成 $g_i\ge0$，乘子的符号也要随之
改变，不能混用两套约定。

KKT 条件通常仍不是充分条件。在凸情形中则有一个简洁的充分性结论。

\begin{theorem}[凸问题中的充分性]
\label{thm:kkt-sufficient-convex}
设 $f,g_i$ 为可微凸函数，$h_j$ 为仿射函数。若可行点 $a$ 与乘子
$\lambda_j\in\R$、$\mu_i\ge0$ 满足
\eqref{eq:kkt-stationarity}--\eqref{eq:kkt-complementarity}，则 $a$ 是全局
极小点。
\end{theorem}

\begin{proof}
对任意可行点 $x$，凸性给出
\[
 f(x)\ge f(a)+\langle\nabla f(a),x-a\rangle.
\]
用驻点方程代换梯度。仿射等式约束使
$\langle\nabla h_j(a),x-a\rangle=0$。又由凸性，
\[
 \langle\nabla g_i(a),x-a\rangle
 \le g_i(x)-g_i(a).
\]
乘以 $-\mu_i$ 后方向反转，并结合 $g_i(x)\le0$、
$\mu_i g_i(a)=0$，得到
$\langle\nabla f(a),x-a\rangle\ge0$。所以 $f(x)\ge f(a)$。
\end{proof}

\begin{example}
在半平面 $x+y\ge1$ 上极小化 $x^2+y^2$。写成
$g(x,y)=1-x-y\le0$。无约束极小点不可行，因此约束活跃。驻点方程
\[
 (2x,2y)+\mu(-1,-1)=0,\qquad x+y=1,\qquad \mu\ge0
\]
给出 $x=y=1/2$、$\mu=1$。目标函数与可行域均凸，故该点是全局极小点。
\end{example}

\section{使用乘子法时的核对顺序}

实际计算中，先确定可行集及活跃约束，再检查正则性；若正则，求解乘子方程
并核实全部原约束。随后依据问题性质选择判定工具：紧集上的直接比较、凸性、
切空间上的二阶型，或局部参数化。若正则性失败，不能简单丢弃该点，而应改换
约束表示、直接参数化，或另作局部分析。

当可行集有端点、尖点、自交或不同维数的分层时，一个统一的光滑等式模型往往
不存在。把每一类边界分别处理，通常比强行套用同一组乘子方程可靠。KKT 理论
把光滑不等式边界纳入了同一计算，但它仍依赖约束资格条件。

\section*{习题}
\addcontentsline{toc}{section}{习题}

\subsection*{A组　等式约束与切空间}
\begin{enumerate}[label=\textbf{\thechapter.\arabic*.}]
 \exerciseitem{factor-functional-through-surjection}
 不使用商空间术语重新证明引理
 \ref{lem:functional-factor-through-surjection}，并说明满射性为何保证乘子唯一。
 \exerciseitem{lagrange-sphere-linear-functional}
 完成球面上线性函数例子的全部候选点与全局比较，包括 $c=0$ 的情形。
 \exerciseitem{lagrange-product-fixed-sum}
 在 $x_1+\cdots+x_n=s$、$x_i>0$ 下求乘积
 $x_1\cdots x_n$ 的最大值，并说明为何不存在最小值。
 \exerciseitem{lagrange-ellipsoid-distance}
 求椭球 $x^2/a^2+y^2/b^2+z^2/c^2=1$ 上到原点距离的最大值和最小值。
 \exerciseitem{two-equality-constraints}
 在 $x+y+z=0$、$x^2+y^2+z^2=1$ 下求 $xyz$ 的极值候选点，并利用
 对称性整理结果。
 \exerciseitem{constraint-representation-degeneracy}
 比较同一集合 $\{0\}$ 的约束 $h(x)=x$ 与 $\widetilde h(x)=x^2$，说明
 乘子条件对约束表示的正则性有何依赖。
 \exerciseitem{tangent-space-intersection-gradients}
 证明多个正则等式约束的切空间是各梯度正交补的交，并计算圆柱
 $x^2+y^2=1$ 与平面 $x+z=0$ 交线在一点的切方向。
\end{enumerate}

\subsection*{B组　二阶判定}
\begin{enumerate}[label=\textbf{\thechapter.\arabic*.},resume]
 \exerciseitem{circle-quadratic-constrained-extrema}
 完整计算正文圆周上 $x^2+2y^2$ 的乘子、切空间和受限二阶型。
 \exerciseitem{sphere-quadratic-eigenvalues}
 设 $A$ 为实对称矩阵。证明二次型 $x^{\mathsf T}Ax$ 在单位球面上的
 驻点正是 $A$ 的单位特征向量，并由此说明极值与最大、最小特征值的关系。
 \exerciseitem{second-variation-curve-derivation}
 从约束曲线二次求导，逐步推导式
 \eqref{eq:constrained-second-variation}。
 \exerciseitem{restricted-hessian-basis-matrix}
 若 $Z$ 的列构成 $T_aM$ 的一组基，证明受限二次型的矩阵为
 $Z^{\mathsf T}\nabla^2_{xx}L(a,\lambda)Z$，并说明如何用它判定正定性。
 \exerciseitem{hessian-objective-alone-fails}
 构造一个约束极小点，使 $\nabla^2 f$ 在切空间外含负方向，但 Lagrange
 Hessian 在切空间上正定；说明检查整个空间会得到错误结论。
 \exerciseitem{constrained-degenerate-second-order}
 在曲线 $y=x^2$ 上研究 $f(x,y)=y$ 和 $f(x,y)=-y$ 在原点的约束极值，
 比较各自的 Lagrange Hessian。
\end{enumerate}

\subsection*{C组　不等式约束}
\begin{enumerate}[label=\textbf{\thechapter.\arabic*.},resume]
 \exerciseitem{kkt-halfspace-quadratic}
 完整核对正文半平面二次优化例子的 KKT 四类条件，并用配方法独立验证答案。
 \exerciseitem{kkt-box-linear-functional}
 在盒子 $0\le x_i\le1$ 上极小化线性函数 $c\cdot x$，写出全部活跃约束
 与乘子，并处理 $c_i=0$ 的非唯一情形。
 \exerciseitem{kkt-disk-linear-objective}
 在闭单位圆盘上极小化 $ax+by$，用 KKT 条件求解，并与
 Cauchy--Schwarz 不等式比较。
 \exerciseitem{kkt-complementarity-sign}
 把约束从 $g_i\le0$ 改写成 $\widehat g_i=-g_i\ge0$，推导相应乘子符号，
 说明互补松弛为何不变。
 \exerciseitem{kkt-qualification-failure}
 对约束 $x^2\le0$ 下极小化 $f(x)=x$，说明极小点存在但 KKT 乘子不存在，
 并指出 LICQ 失败之处。
 \exerciseitem{kkt-not-sufficient-nonconvex}
 给出满足 KKT 条件却不是局部极小点的光滑不等式约束例子。
 \exerciseitem{convex-kkt-sufficiency-detail}
 补全定理 \ref{thm:kkt-sufficient-convex} 证明中的每个不等号，并说明
 若 $h_j$ 不是仿射函数，证明在哪一步失效。
 \optionalexerciseitem{kkt-coordinate-proof-reconstruction}
 详细构造 KKT 定理证明中的补充坐标 $z$，并用逆函数定理写出局部可行集
 的坐标形式。
\end{enumerate}
